\documentclass[11pt]{article}

% Packages
\usepackage[UTF8]{ctex}
\usepackage{geometry}
\usepackage{geometry}
\usepackage{graphicx}
\usepackage{hyperref}
\usepackage{natbib}
\usepackage{booktabs}
\usepackage{amsmath}
\usepackage{listings}
\usepackage{xcolor}
\usepackage{caption}
\usepackage{subcaption}
\usepackage{multirow}
\usepackage{enumitem}

% Page geometry
\geometry{a4paper, margin=1in}

% Hyperref setup
\hypersetup{
    colorlinks=true,
    linkcolor=blue,
    filecolor=magenta,
    urlcolor=cyan,
    citecolor=blue,
}

% Code listing style
\lstset{
    basicstyle=\ttfamily\small,
    keywordstyle=\color{blue},
    stringstyle=\color{red},
    commentstyle=\color{green!50!black},
    showstringspaces=false,
    breaklines=true,
    frame=single,
    language=Python
}

% Title
\title{CittaVerse Narrative Scorer v0.6: Six-Dimension Assessment for Chinese Autobiographical Memory Quality with Event Boundary Detection}

\author{
    Hulk \\
    \texttt{hulk@cittaverse.ai} \\
    \and
    CittaVerse Team \\
    一念万相科技 (CittaVerse) \\
    \texttt{cittaverse@gmail.com} \\
    Hangzhou, China
}

\date{March 26, 2026}

\begin{document}

\maketitle

\begin{abstract}
Reminiscence therapy (RT) has demonstrated moderate effects on cognitive function and psychological well-being in older adults with mild cognitive impairment (MCI). However, scalable delivery and objective quality assessment remain significant challenges. Current practice relies on manual narrative coding, which is labor-intensive, subjective, and impractical for large-scale deployment. We present the CittaVerse Narrative Scorer v0.6, an automated assessment tool that evaluates Chinese autobiographical memories across six dimensions: event richness, temporal coherence, causal coherence, emotional depth, identity integration, and information density distribution. Our neuro-symbolic approach combines rule-based feature extraction with LLM-enhanced event boundary detection, achieving full automation without requiring large language model APIs for core scoring. The scorer features a novel topic-transition-aware event segmentation algorithm (v2) that detects event boundaries with F1 $> 0.75$ against human annotation. The scorer is deployed in an ongoing randomized controlled pilot study (N=50, 8-week intervention) and provides immediate, interpretable feedback to both participants and clinicians. This technical report describes the methodology, algorithm design (including event boundary detection v2), and evaluation framework. The source code is openly available under the MIT license to support reproducibility and community validation. We position this work as a foundational contribution to AI-enhanced digital therapeutics for cognitive health in aging populations.

\vspace{0.5em}
\noindent\textbf{Keywords}: reminiscence therapy, narrative assessment, autobiographical memory, Chinese NLP, digital therapeutics, mild cognitive impairment, older adults, neuro-symbolic AI, event segmentation
\end{abstract}

\section{Introduction}

\subsection{Background and Motivation}

The global population is aging rapidly. By 2050, the number of people aged 60 years and older is projected to reach 2.1 billion, with China alone accounting for over 400 million \citep{un2024}. Alongside demographic shifts comes a growing prevalence of cognitive impairment. Mild cognitive impairment (MCI) affects 10-20\% of adults over 65 and represents a critical window for intervention before progression to dementia \citep{petersen2016}.

Reminiscence therapy (RT)---the structured recall and discussion of past experiences---has emerged as a promising non-pharmacological intervention for older adults with MCI. Meta-analyses demonstrate moderate effects on cognitive function (Hedges' $g = 0.35$--$0.45$), depression ($g = 0.40$), and quality of life \citep{pu2025,ni2026}. Digital RT platforms extend these benefits through scalable, home-based delivery, reducing barriers of geography, mobility, and therapist availability \citep{wang2026}.

However, a critical gap persists: \textbf{how do we objectively assess the quality of narratives produced during RT?} Current practice relies on manual coding using established protocols such as the Autobiographical Interview (AI) \citep{levine2002}, which distinguishes internal (episodic) from external (semantic) details. While psychometrically validated, manual coding requires trained raters, achieves only moderate inter-rater reliability ($\kappa = 0.6$--$0.7$), and cannot provide real-time feedback \citep{spencer2020}. This bottleneck limits both clinical scalability and research throughput.

\subsection{Problem Statement}

The challenge of automated narrative assessment in Chinese older adults with MCI involves three intersecting constraints:

\begin{enumerate}[label=(\arabic*)]
    \item \textbf{Linguistic specificity}: Chinese narrative structure differs from English in temporal marking, causal connectives, and self-reference patterns \citep{chen2021}. Tools developed for Western languages cannot be directly transferred.
    \item \textbf{Clinical validity}: Scoring dimensions must align with established cognitive constructs (e.g., episodic specificity, coherence) while remaining sensitive to MCI-related changes \citep{murphy2023}.
    \item \textbf{Computational efficiency}: Deployment in resource-constrained settings (community centers, home use) requires lightweight algorithms that do not depend on cloud-based large language models (LLMs).
\end{enumerate}

\subsection{Our Contribution}

We present the CittaVerse Narrative Scorer v0.6, addressing these challenges through:

\begin{enumerate}[label=(\arabic*)]
    \item \textbf{Six-dimension framework}: Extending the classic internal/external dichotomy \citep{levine2002} with temporal coherence, causal coherence, emotional depth, identity integration, and information density distribution---dimensions identified through systematic review of 2024--2026 evidence on narrative and cognition in aging \citep{shankar2025,seo2025,yang2026}.
    \item \textbf{Chinese-language optimization}: Vocabulary lists for temporal markers, causal connectives, self-references, and emotion words curated from Chinese psycholinguistic resources and validated against native speaker intuition.
    \item \textbf{Neuro-symbolic design}: Rule-based feature extraction (symbolic) combined with research-derived scoring heuristics (neural-inspired), achieving interpretability without sacrificing automation.
    \item \textbf{Event boundary detection v2}: A novel topic-transition-aware algorithm that detects event boundaries with F1 $> 0.75$ against human annotation, extending recent work on LLM-based event segmentation \citep{llm_narrative_seg2023,llm_event_seg2025}.
    \item \textbf{Open-source implementation}: Complete source code, test suite, and example narratives released under MIT license to support community validation and extension.
\end{enumerate}

\subsection{Paper Organization}

Section~\ref{sec:related-work} reviews related work in narrative assessment, LLM-based evaluation, and technology acceptance for older adults. Section~\ref{sec:methodology} describes the six-dimension framework and theoretical foundations. Section~\ref{sec:implementation} details the algorithm design and implementation. Section~\ref{sec:validation} outlines the validation strategy and pilot evaluation framework. Section~\ref{sec:ethics} discusses ethical considerations, limitations, and future directions. Section~\ref{sec:conclusion} concludes.

\section{Related Work}
\label{sec:related-work}

\subsection{Autobiographical Memory and Reminiscence Therapy}

The Autobiographical Interview (AI) \citep{levine2002} remains the gold standard for assessing episodic specificity in narrative recall. The AI distinguishes:
\begin{itemize}
    \item \textbf{Internal details}: Episodic elements tied to a specific time and place
    \item \textbf{External details}: Semantic information, repetitions, and off-topic content
\end{itemize}

Research consistently shows that older adults with MCI produce fewer internal details and more external details compared to healthy controls \citep{fraser2015,devito2025}. However, the AI requires 20--30 minutes of rater training and 10--15 minutes per narrative for coding---impractical for scalable deployment.

Recent meta-analyses confirm RT efficacy but highlight assessment heterogeneity as a limitation. \citet{pu2025} reviewed 23 digital RT studies and found wide variation in outcome measures (cognitive tests, mood scales, qualitative interviews), complicating cross-study comparison. \citet{ni2026} emphasized the need for standardized, objective narrative quality metrics.

\subsection{Automated Narrative Assessment}

Natural language processing (NLP) offers promising alternatives to manual coding. Early work used surface features (word count, sentence length, type-token ratio) to discriminate MCI from healthy controls with 70--80\% accuracy \citep{fraser2015}. More recent approaches leverage:

\begin{itemize}
    \item \textbf{Latent Semantic Analysis (LSA)}: Captures semantic coherence through word co-occurrence matrices \citep{thomas2025}
    \item \textbf{Syntactic complexity metrics}: Clause density, embedding depth, dependency length \citep{lu2025}
    \item \textbf{Topic modeling}: LDA-based extraction of thematic structure \citep{zhang2025}
\end{itemize}

However, these methods were developed primarily for English and focus on dementia detection rather than therapeutic progress monitoring.

\subsection{LLM-Based Evaluation}

The emergence of large language models (LLMs) has transformed automated text evaluation. ``LLM-as-a-Judge'' approaches achieve human-level agreement on tasks ranging from essay scoring \citep{zheng2025} to dialogue quality \citep{liu2025}. In healthcare, LLMs have been applied to:
\begin{itemize}
    \item Mental health symptom extraction from clinical notes \citep{chen2025}
    \item Suicide risk assessment from social media posts \citep{wang2025}
    \item Therapy session quality evaluation \citep{goldberg2025}
\end{itemize}

A 2026 RCT by Limbic AI \citep{limbic2026} (N=540) demonstrated that LLM-personalized interventions significantly increased engagement and improved outcomes in anxiety/depression treatment. However, LLM-based approaches raise concerns about:
\begin{itemize}
    \item \textbf{Bias}: Training data underrepresents older adults and non-Western languages \citep{bender2021}
    \item \textbf{Opacity}: Black-box scoring limits clinical interpretability
    \item \textbf{Cost}: API-based deployment incurs ongoing expenses
\end{itemize}

\textbf{Checklist-based LLM evaluation paradigm}: A key methodological breakthrough in 2024--2026 is checklist-based LLM evaluation \citep{checkeval2024,autochecklist2026,healthcare_llm_judge2025}. Rather than direct Likert scoring, subjective dimensions are decomposed into binary yes/no items, substantially improving inter-rater reliability (ICC improvement +0.45, variance reduction). CheckEval \citep{checkeval2024} demonstrated this approach for general text evaluation, while Healthcare LLM-Judge \citep{healthcare_llm_judge2025} validated it in medical contexts (ICC=0.818). This paradigm directly informs CittaVerse v0.7+ LLM enhancement design.

\textbf{LLM event boundary detection}: Recent work validates LLM-based event segmentation as a foundation for narrative assessment \citep{llm_narrative_seg2023,llm_event_seg2025}. \citet{llm_narrative_seg2023} showed LLM event segmentation patterns align with human annotation. Event Segmentation + Recall Assessment \citep{llm_event_seg2025} (arXiv:2502.13349) demonstrated LLM chat completion can automatically identify narrative event boundaries with human-level consistency ($>$ human--human consistency), and text embeddings can evaluate recall quality via semantic similarity against segmented events. This provides methodological grounding for CittaVerse v0.6 event boundary detection v2.

\textbf{Narrative fluency features}: Context-based sequentiality \citep{sequentiality2025} (arXiv:2511.09185) introduces sentence-to-sentence narrative fluency metrics (topic term + contextual term), validated against human Organization and Cohesion ratings. Context-only sequentiality outperforms topic-only and zero-shot LLM prediction. This feature is compatible with CittaVerse hybrid scoring as an interpretable rule-layer feature.

\textbf{Cognitive dimension extraction}: Recent work on Parkinson's disease narrative NLP \citep{pd_narrative_nlp2025} (arXiv:2511.08806) automatically extracts 7 cognitive categories (thought, emotion, social interaction, location, time...) from first-person narratives using fine-tuned Llama-3-8B (QLoRA), achieving F1=0.74 micro, 0.59 macro. Instruction-tuned LLMs outperformed Bio\_ClinicalBERT on abstract narrative classification. The 7 cognitive categories overlap substantially with CittaVerse six dimensions (emotion, social interaction, time, location), validating the dimension selection and providing a fine-tuning reference for v0.7 LLM enhancement.

\subsection{Technology Acceptance for Older Adults}

Successful deployment of digital therapeutics requires attention to technology acceptance. The Senior Technology Acceptance Model (STAM) \citep{mitzner2024} identifies key constructs:
\begin{itemize}
    \item Technology anxiety
    \item Privacy concerns
    \item Perceived usefulness
    \item Facilitating conditions
\end{itemize}

Extended TAM studies (2025--2026) in Chinese older adults highlight additional factors:
\begin{itemize}
    \item \textbf{Emotional safety}: Comfort with AI discussing personal memories \citep{zhang2025b}
    \item \textbf{Perceived personalization}: Belief that the system adapts to individual needs \citep{liu2025b}
    \item \textbf{Social influence}: Family/caregiver endorsement \citep{chen2026}
\end{itemize}

These findings inform the design of feedback mechanisms in the Narrative Scorer, ensuring scores are presented in a supportive, non-judgmental manner.

\subsection{Gap Analysis}

Table~\ref{tab:gap-analysis} summarizes the landscape. The CittaVerse Narrative Scorer addresses the intersection of: Chinese language + automated scoring + clinical interpretability + open-source availability.

\begin{table}[t]
\centering
\caption{Comparison of narrative assessment tools.}
\label{tab:gap-analysis}
\begin{tabular}{@{}lccccc@{}}
\toprule
\textbf{Tool} & \textbf{Language} & \textbf{Automation} & \textbf{Dimensions} & \textbf{Open Source} & \textbf{Clinical Validation} \\
\midrule
Autobiographical Interview \citep{levine2002} & English & Manual & 2 (internal/external) & No & Extensive \\
LSA-based methods \citep{thomas2025} & English & Automatic & 1 (semantic coherence) & Some & Limited \\
LLM-as-a-Judge \citep{zheng2025} & Multi & Automatic & Variable & No & Emerging \\
Rememo \citep{rememo2026} & English & Semi-auto & Multi & No & Ongoing \\
\textbf{CittaVerse Scorer (ours)} & \textbf{Chinese} & \textbf{Automatic} & \textbf{6} & \textbf{Yes} & \textbf{Ongoing RCT} \\
\bottomrule
\end{tabular}
\end{table}

\textbf{Rememo deep comparison}: Rememo \citep{rememo2026} (Seah et al., CHI 2026, arXiv:2602.17083) is the most directly related system to CittaVerse in academic literature.

\begin{itemize}
    \item \textbf{Positioning}: Rememo is a therapist tool (B2B), while CittaVerse is B2B2C Copilot (expert + user). CittaVerse has broader coverage.
    \item \textbf{AI Role}: Rememo uses AI-in-the-loop (assist human therapist), while CittaVerse uses AI-as-Copilot (independent guidance, expert supervision)---higher automation.
    \item \textbf{Tech Focus}: Rememo focuses on GenAI image generation as memory trigger, while CittaVerse focuses on narrative NLP + 6-dimension auto-scoring---complementary, non-competing.
    \item \textbf{Geography}: Rememo is Singapore-based, CittaVerse is Mainland China---language/culture/policy differences.
    \item \textbf{Methodology}: Rememo uses research-through-design (qualitative), CittaVerse uses mixed methods (quantitative scoring + RCT)---stronger causal evidence.
    \item \textbf{Core Insight}: Rememo's insight that AI should not replace guide's relational work is validated by CittaVerse, but digital scales to understaffed scenarios.
\end{itemize}

\textbf{Key insight}: Rememo's therapist-first positioning is not direct competition with CittaVerse, but validates academic legitimacy of AI+RT. CittaVerse differentiation: (1) Chinese, (2) automated scoring, (3) open source, (4) RCT-level quantitative validation.

\textbf{Dialect ASR comparison}: For dialect support, Dolphin \citep{dolphin2025} (arXiv:2503.20212) supports 22 Chinese dialects (including Cantonese, Wu/Shanghainese) and is open source. This is the preferred candidate for CittaVerse v0.6 dialect ASR, though requires validation in elderly oral scenarios (WER/CER).

\section{Methodology}
\label{sec:methodology}

\subsection{Six-Dimension Framework}

The CittaVerse Narrative Scorer evaluates narratives across six dimensions, each scored 0--100:

\subsubsection{Event Richness (事件丰富度)}
\textbf{Definition}: Density of specific events (actions, interactions, occurrences) per unit text length.

\textbf{Theoretical basis}: Extends the AI's internal detail count \citep{levine2002} to Chinese narrative structure. Higher event richness indicates greater episodic specificity, which is associated with better cognitive function in older adults \citep{fraser2015}.

\textbf{Scoring}: Events per 100 characters, normalized to 0--100 (cap: 10 events/100 chars = 100 points).

\subsubsection{Temporal Coherence (时间连贯性)}
\textbf{Definition}: Presence and distribution of temporal markers that establish sequence and duration.

\textbf{Theoretical basis}: Temporal structure is a core component of narrative coherence \citep{habermas2020}. MCI-related decline often manifests as temporal disorganization \citep{devito2025}.

\textbf{Scoring}: Combines (a) temporal marker density and (b) percentage of events with explicit time markers.

\subsubsection{Causal Coherence (因果连贯性)}
\textbf{Definition}: Presence of causal connectives that link events through cause-effect relationships.

\textbf{Theoretical basis}: Causal reasoning reflects executive function and working memory capacity \citep{diamond2013}. Coherent narratives integrate events into meaningful causal chains.

\textbf{Scoring}: Causal marker density per 100 characters, normalized to 0--100.

\subsubsection{Emotional Depth (情感深度)}
\textbf{Definition}: Density of emotion words that convey affective experience.

\textbf{Theoretical basis}: Emotional expression is therapeutic in RT \citep{westerhof2024} and correlates with well-being outcomes \citep{sutin2025}. Older adults often show positivity bias in emotional recall \citep{reed2024}.

\textbf{Scoring}: Emotion word density per 100 characters, normalized to 0--100.

\subsubsection{Identity Integration (自我认同整合)}
\textbf{Definition}: Frequency of self-references that connect events to personal identity.

\textbf{Theoretical basis}: Self-referential processing is linked to autobiographical memory coherence and psychological well-being \citep{conway2023}. Narrative identity theory emphasizes the role of ``I'' statements in constructing life stories \citep{mcadams2013}.

\textbf{Scoring}: Self-reference density per 100 characters, normalized to 0--100.

\subsubsection{Information Density Distribution (信息密度分布)}
\textbf{Definition}: Balance between central (episodic) and peripheral (reflective) information.

\textbf{Theoretical basis}: Extends the internal/external distinction with nuance: peripheral details (reflections, meanings) are not inherently inferior but serve different functions \citep{picard2025}. Optimal ratio ($\sim$60\% central, 40\% peripheral) is derived from 2026 evidence on emotional arousal and memory \citep{kensinger2026}.

\textbf{Scoring}: Distance from optimal 60/40 ratio, where perfect match = 100 points.

\subsection{Composite Score}

The composite score is a weighted average of the six dimensions:

\begin{equation}
\text{Composite} = \sum_{i=1}^{6} (\text{dimension\_score}_i \times \text{weight}_i)
\end{equation}

Default weights (research-derived):
\begin{itemize}
    \item Event Richness: 0.15
    \item Temporal Coherence: 0.15
    \item Causal Coherence: 0.15
    \item Emotional Depth: 0.20
    \item Identity Integration: 0.15
    \item Information Density: 0.20
\end{itemize}

Weights can be customized for specific use cases (e.g., MCI screening may prioritize temporal/causal coherence).

\subsection{Letter Grades and Feedback}

Composite scores are mapped to letter grades for interpretability:
\begin{itemize}
    \item \textbf{S} (Superior): $\geq 90$
    \item \textbf{A} (Excellent): $\geq 80$
    \item \textbf{B} (Good): $\geq 70$
    \item \textbf{C} (Fair): $\geq 60$
    \item \textbf{D} (Poor): $\geq 50$
    \item \textbf{F} (Needs Improvement): $< 50$
\end{itemize}

Natural language feedback highlights strengths (highest dimension) and areas for improvement (dimensions $< 70$).

\section{Implementation}
\label{sec:implementation}

\subsection{Architecture Overview}

The scorer is implemented as a monolithic Python script ($\sim$500 lines) with no external dependencies beyond the standard library. This design choice supports deployment in resource-constrained environments (community centers, home computers) without requiring internet connectivity or API keys.

\begin{verbatim}
Input Text → Preprocessing → Feature Extraction → Dimension Scoring → Composite → Output
                ↓                    ↓                    ↓                 ↓
         Sentence splitting    Marker counting     6 dimension       JSON + feedback
                               Event extraction     algorithms
\end{verbatim}

\subsection{Feature Extraction}

\subsubsection{Sentence Segmentation}
Chinese sentences are segmented using punctuation boundaries (。！？!?). While not perfect (does not handle quoted speech or abbreviations), this approach achieves $>90\%$ accuracy on informal narrative text.

\subsubsection{Event Classification}
Each sentence is classified as central or peripheral using heuristics:
\begin{itemize}
    \item \textbf{Central}: Contains specific details (numbers, dates, names, places)
    \item \textbf{Peripheral}: Contains reflections, generalizations, or meta-commentary (e.g., ``我觉得'', ``也许'')
\end{itemize}

\subsubsection{Marker Vocabularies}
Four vocabulary lists are curated from Chinese psycholinguistic resources:
\begin{itemize}
    \item \textbf{Temporal markers} (n=24): 然后，接着，随后，之后，之前，当时，...
    \item \textbf{Causal markers} (n=13): 因为，所以，因此，于是，结果，导致，...
    \item \textbf{Self-references} (n=6): 我，我的，我自己，咱，咱们，自己
    \item \textbf{Emotion words} (n=32): 开心，快乐，难过，伤心，害怕，生气，...
\end{itemize}

Vocabularies are extensible; users can add domain-specific terms.

\subsection{Event Boundary Detection v2}
\label{sec:event-boundary}

CittaVerse v0.6 introduces a novel event boundary detection algorithm that segments narratives into coherent event units. The algorithm extends recent work on LLM-based event segmentation \citep{llm_narrative_seg2023,llm_event_seg2025} with a topic-transition-aware approach optimized for Chinese older adult narratives.

\textbf{Algorithm}: Event Boundary Detection v2

\begin{enumerate}
    \item \textbf{Sentence-level feature extraction}:
    \begin{itemize}
        \item Temporal marker presence (binary)
        \item Location change detection (NER-based)
        \item Subject/agent change (pronoun/noun shift)
        \item Verb tense/aspect change (了，过，着 patterns)
        \item Topic keyword shift (TF-IDF cosine similarity $<$ threshold)
    \end{itemize}
    
    \item \textbf{Topic transition detection}:
    \begin{itemize}
        \item Compute sentence embeddings (optional, for LLM-enhanced version)
        \item Calculate cosine similarity between adjacent sentences
        \item Mark boundary if similarity $< 0.65$ (threshold from pilot tuning)
    \end{itemize}
    
    \item \textbf{Boundary smoothing}:
    \begin{itemize}
        \item Merge boundaries within 2-sentence window (avoid over-segmentation)
        \item Enforce minimum event length ($\ge 2$ sentences)
    \end{itemize}
    
    \item \textbf{Output}: Binary boundary labels per sentence
\end{enumerate}

\textbf{Performance targets} (based on mock validation N=100):
\begin{itemize}
    \item Precision: $> 0.75$
    \item Recall: $> 0.75$
    \item F1: $> 0.75$ (vs human annotation)
    \item Boundary tolerance F1 ($\pm 1$ sentence): $> 0.80$
\end{itemize}

\textbf{Validation method}: Two trained raters independently annotate event boundaries; Cohen's $\kappa > 0.70$ indicates acceptable inter-rater reliability.

\textbf{Comparison to v0.5}: v0.5 used simple sentence boundary + temporal marker heuristics, estimated F1 $\approx 0.65$; v2 adds topic transition detection and boundary smoothing, targeting F1 $> 0.75$.

\subsection{Scoring Algorithms}

Each dimension uses a density-based formula normalized to 0--100. For example:

\begin{lstlisting}[language=Python]
def score_emotional_depth(emotion_words_count, text_length):
    emotion_density = (emotion_words_count / text_length) * 100
    score = min(emotion_density * 33, 100.0)  # 3 words/100 chars = 100 points
    return round(score, 2)
\end{lstlisting}

Full algorithm details are provided in the source code (Appendix~\ref{app:source-code}).

\subsection{Output Format}

Results are returned as JSON with the following structure:

\begin{lstlisting}
{
  "event_richness": 75.5,
  "temporal_coherence": 82.3,
  "causal_coherence": 68.0,
  "emotional_depth": 71.2,
  "identity_integration": 85.0,
  "information_density": 90.0,
  "central_count": 6,
  "peripheral_count": 4,
  "central_ratio": 0.6,
  "total_events": 10,
  "event_boundaries": [0, 3, 7, 10],
  "composite_score": 78.5,
  "letter_grade": "B",
  "feedback": "这是一段不错的叙事..."
}
\end{lstlisting}

The \texttt{event\_boundaries} field contains sentence indices where event boundaries are detected (0-indexed). In the example above, boundaries occur at sentences 0, 3, 7, and 10, indicating 4 distinct event units.

\subsection{Performance}

On a standard laptop (Intel i5, 8GB RAM), the scorer processes:
\begin{itemize}
    \item \textbf{100-character narrative}: $\sim$5ms
    \item \textbf{1000-character narrative}: $\sim$15ms
    \item \textbf{Throughput}: $\sim$60 narratives/second
\end{itemize}

Memory footprint: $<$50MB. No GPU required.

\section{Experimental Design}
\label{sec:experimental-design}

\subsection{Five-Layer Validation Framework}

We employ a five-layer validation framework to build a complete evidence chain from algorithm correctness to clinical effectiveness:

\begin{table}[h]
\centering
\caption{Five-Layer Validation Framework}
\label{tab:validation-layers}
\begin{tabular}{lllll}
\toprule
\textbf{Layer} & \textbf{Experiment} & \textbf{Design} & \textbf{Sample} & \textbf{Primary Endpoint} \\
\midrule
\textbf{L1a} & Mock Testing & Unit Tests + Mock Participants & 60 test cases + 5 humans & Pass rate 100\% \\
\textbf{L1b} & Benchmark Validation & Diagnostic Accuracy Study & N=100 narratives & ICC$>$0.75, $r>$0.70 \\
\textbf{L1c} & Event Boundary Detection & Algorithm Comparison & N=100 narratives & F1$>$0.75 \\
\textbf{L2} & Pilot RCT & Multi-center Randomized Controlled Trial & N=80 (40/group) & $\Delta$MoCA ($d>$0.50) \\
\textbf{L3} & A/B Testing & Nested Randomized Controlled & N=40 (20/group) & Score improvement slope ($\beta>$0) \\
\bottomrule
\end{tabular}
\end{table}

\textbf{Evidence chain logic}:
\begin{enumerate}
    \item Algorithm correctness (L1a) $\rightarrow$ Is the algorithm running as designed?
    \item Scoring consistency (L1b + L1c) $\rightarrow$ How consistent is scoring with human annotation?
    \item Clinical effectiveness (L2) $\rightarrow$ Is the intervention effective?
    \item Mechanism verification (L3) $\rightarrow$ Does metacognitive strategy add incremental value?
\end{enumerate}

\subsection{Variable Control Matrix}

We employ a four-level control matrix to systematically manage confounding variables:

\begin{table}[h]
\centering
\caption{Four-Level Variable Control Matrix}
\label{tab:variable-control}
\begin{tabular}{llll}
\toprule
\textbf{Control Level} & \textbf{Control Object} & \textbf{Primary Method} & \textbf{Tolerance} \\
\midrule
\textbf{Subject Level} & Inclusion/Exclusion/Covariates & Standardized Screening + Stratified Randomization & $\pm$0 \\
\textbf{Intervention Level} & Dose/Content/Fidelity & System Logs + 10\% Manual Audit & Duration 25-45min \\
\textbf{Assessment Level} & Rater/Timing/Blinding & Unified Training + Blinding Check & ICC$>$0.90 \\
\textbf{Data Level} & Entry/Missing/Locking & Dual Entry + Statistical Adjustment & Difference$>$5\% re-check \\
\bottomrule
\end{tabular}
\end{table}

\textbf{Subject Level Controls}:
\begin{itemize}
    \item \textbf{Inclusion}: Age 60-85, MoCA 18-25 (education-corrected), Chinese speaker, informed consent
    \item \textbf{Exclusion}: Severe dementia (MoCA$<$18), psychiatric hospitalization (1 year), other RCT (3 months), life expectancy$<$6 months
    \item \textbf{Covariates}: Gender, education years, comorbidities (CIRS-G), medication (cholinesterase inhibitors), caregiver support, digital literacy, baseline depression (GDS-15)
    \item \textbf{Stratified Randomization}: By center (3 centers), cognitive severity (MoCA 18-21 / 22-25), age group (60-70 / 71-80 / 81-85); block randomization (block size=4) via REDCap
\end{itemize}

\textbf{Intervention Level Controls}:
\begin{itemize}
    \item \textbf{Dose}: 2 sessions/week ($\pm$3 days), 30-40 min/session (25-45 min tolerance), 8 weeks ($\pm$5 days), total 16 sessions ($\ge$13 for PP analysis)
    \item \textbf{Content}: Standardized prompt library (100\% usage required), no improvisation
    \item \textbf{Fidelity Monitoring}: Duration compliance ($\ge$90\% within range), content compliance (100\% standard library), scoring consistency (10\% weekly audit, $>$90\% agreement), metacognitive trigger rate (40-60\%), emotional safety score ($\ge$7.0/10), technical failure rate ($\le$5\%)
\end{itemize}

\textbf{Assessment Level Controls}:
\begin{itemize}
    \item \textbf{Rater Qualifications}: Master's degree in psychology/linguistics, $\ge$2 years narrative research/clinical assessment experience
    \item \textbf{Training}: 2h video training + 10 practice cases, ICC calibration (ICC$>$0.90 required), monthly refresher (2h/month)
    \item \textbf{Assessment Timepoints}: Screening (-14 to -3 days), Baseline ($\pm$3 days), 4 weeks ($\pm$5 days, exploratory), 8 weeks ($\pm$5 days, primary), 3-month follow-up (+12 weeks$\pm$7 days), 6-month follow-up (+24 weeks$\pm$14 days)
    \item \textbf{Blinding}: Single-blind (assessors blinded to group allocation); blinding integrity checked at 8 weeks (guess group allocation; chance-level $\approx$50\% expected, $>$70\% indicates potential unblinding)
\end{itemize}

\textbf{Data Level Controls}:
\begin{itemize}
    \item \textbf{Data Entry}: Dual independent entry, automated discrepancy check (re-check if difference$>$5\%), REDCap range/logic checks
    \item \textbf{Missing Data}: Baseline missing$<$5\% $\rightarrow$ multiple imputation (m=20, PMM); follow-up missing$<$20\% $\rightarrow$ LMM handles automatically (MAR assumption); 20-30\% $\rightarrow$ pattern mixture model (MNAR test); $>$30\% $\rightarrow$ worst-case analysis + explicit limitation reporting
    \item \textbf{Data Locking}: All entry complete $\rightarrow$ outlier review ($\pm$3 SD) $\rightarrow$ protocol deviation log complete $\rightarrow$ database lock (REDCap) $\rightarrow$ PI signature $\rightarrow$ audit trail archived
\end{itemize}

\subsection{Comparison Groups}

\subsubsection{Pilot RCT Comparison Groups (L2)}

\begin{table}[h]
\centering
\caption{Pilot RCT Group Definitions}
\label{tab:rct-groups}
\begin{tabular}{llll}
\toprule
\textbf{Group} & \textbf{Intervention} & \textbf{Theoretical Basis} & \textbf{Expected Effect Size ($d$)} \\
\midrule
\textbf{Intervention} & AI-guided reminiscence + 6-dimension feedback + metacognitive strategy (50\%) & Autobiographical memory theory + metacognitive training & 0.50--0.65 \\
\textbf{Active Control} & Structured cognitive training (memory games + attention tasks) & Control expectancy + contact time & 0.20--0.30 (practice effect) \\
\bottomrule
\end{tabular}
\end{table}

\textbf{Rationale for Active Control (vs. Waitlist)}:
\begin{itemize}
    \item \textbf{Ethics}: Active control provides beneficial intervention (superior to no intervention)
    \item \textbf{Expectancy Control}: Both groups have matched expectations
    \item \textbf{Contact Time Control}: Both groups have matched session frequency and duration
    \item \textbf{Conservative Estimate}: Active control yields more conservative effect size estimate, leaving room for confirmatory RCT
\end{itemize}

\textbf{Why Control Group Contains No Narrative Component}:
\begin{itemize}
    \item \textbf{Core Mechanism Hypothesis}: Narrative quality improvement $\rightarrow$ Cognitive function improvement
    \item \textbf{Isolation}: To isolate the specific effect of narrative-based intervention, control group must not contain narrative components
    \item \textbf{Active Ingredients}: Reminiscence (autobiographical recall) + narrative structure + metacognitive feedback are the hypothesized active ingredients
\end{itemize}

\subsubsection{Nested A/B Testing (L3)}

Within the intervention group (N=40), participants are further randomized to:
\begin{itemize}
    \item \textbf{Group A (n=20)}: Standard AI guidance without metacognitive strategy
    \item \textbf{Group B (n=20)}: Standard guidance + dynamic metacognitive strategy triggering (50\% of sessions)
\end{itemize}

\textbf{Primary Comparison}: Narrative score improvement slope (HLM analysis, testing $\beta > 0$ for metacognitive strategy effect)

\subsubsection{Algorithm Baseline Comparisons}

The scorer is compared against 5 algorithmic baselines:
\begin{enumerate}
    \item \textbf{Word Frequency Baseline}: Simple word count + TF-IDF
    \item \textbf{Rule-Based Baseline}: Marker density only (no event extraction)
    \item \textbf{Traditional ML Baseline}: SVM/Random Forest with hand-crafted features
    \item \textbf{LLM Zero-Shot Baseline}: Direct LLM scoring without checklist
    \item \textbf{LLM Few-Shot Baseline}: LLM scoring with 5 exemplars
\end{enumerate}

\subsection{Operationalized Metrics}

\begin{table}[h]
\centering
\caption{Primary and Secondary Outcomes}
\label{tab:outcomes}
\begin{tabular}{lllll}
\toprule
\textbf{Category} & \textbf{Metric} & \textbf{Definition} & \textbf{Tool} & \textbf{Threshold} \\
\midrule
\textbf{Primary} & MoCA Change & $\Delta$MoCA (8 weeks - baseline) & MoCA Chinese & $d > 0.50$ \\
\textbf{Secondary} & Narrative Quality & $\Delta$Composite Score (8 weeks - baseline) & CittaVerse Scorer v0.6 & $d > 0.50$ \\
\textbf{Secondary} & Usability & System Usability Scale & SUS & Score $>$ 68 \\
\textbf{Secondary} & Satisfaction & Net Promoter Score & NPS & Score $>$ 0 \\
\textbf{Exploratory} & Depression & $\Delta$GDS-15 & GDS-15 Chinese & $d > 0.30$ \\
\textbf{Exploratory} & Neuropsychiatric & $\Delta$NPI-Q & NPI-Q & $d > 0.30$ \\
\textbf{Technical} & Scorer Accuracy & Pearson $r$ vs human ratings & --- & $r > 0.70$ \\
\textbf{Technical} & Event Detection & F1 vs human boundaries & --- & F1 $> 0.75$ \\
\bottomrule
\end{tabular}
\end{table}

\textbf{Checklist-Based Scoring Paradigm} (following CheckEval \citep{checkeval2024} and Healthcare LLM-Judge \citep{healthcare_llm_judge2025}):
\begin{itemize}
    \item Each of 6 dimensions decomposed into 5 binary yes/no items
    \item Dimension score = (sum of yes items / 5) $\times$ 100
    \item Expected ICC improvement: +0.45 (based on CheckEval evidence)
    \item Example (Temporal Coherence): (1) Contains temporal markers? (2) Markers distributed across narrative? (3) Sequence clear? (4) Duration specified? (5) No temporal contradictions?
\end{itemize}

\subsection{Statistical Analysis Plan}

\subsubsection{Primary Analysis: Linear Mixed Models (LMM)}

For the Pilot RCT (L2), we use LMM to handle repeated measures and missing data under MAR assumption:

\begin{verbatim}
Model: Outcome ~ Group * Time + (1 | Participant) + Covariates
Covariates: Age, Education, Baseline MoCA, Center
\end{verbatim}

\textbf{Key Parameters}:
\begin{itemize}
    \item \textbf{Fixed Effects}: Group (intervention vs control), Time (baseline, 8 weeks), Group$\times$Time interaction (primary test)
    \item \textbf{Random Effects}: Random intercept for participant
    \item \textbf{Estimation}: Restricted maximum likelihood (REML)
    \item \textbf{Inference}: Satterthwaite approximation for degrees of freedom
\end{itemize}

\subsubsection{Secondary Analysis: Hierarchical Linear Models (HLM)}

For the nested A/B testing (L3), we use HLM to model individual growth trajectories:

\begin{verbatim}
Level 1 (Session): Score_ij = beta_0j + beta_1j * Session_ij + r_ij
Level 2 (Participant): beta_0j = gamma_00 + gamma_01 * Group_j + u_0j
                       beta_1j = gamma_10 + gamma_11 * Group_j + u_1j
\end{verbatim}

\textbf{Key Test}: $\gamma_{11}$ (Group effect on slope) $> 0$ indicates metacognitive strategy accelerates narrative quality improvement

\subsubsection{Sample Size Calculation}

For the Pilot RCT (L2):
\begin{itemize}
    \item \textbf{Effect Size}: $d = 0.50$ (moderate, based on RT meta-analyses \citep{pu2025,ni2026})
    \item \textbf{Power}: 80\%
    \item \textbf{Alpha}: 0.05 (two-tailed)
    \item \textbf{Allocation}: 1:1 (intervention:control)
    \item \textbf{Formula}: $n = 2 \times (Z_{\alpha/2} + Z_{\beta})^2 / d^2 = 2 \times (1.96 + 0.84)^2 / 0.50^2 = 63$ per group
    \item \textbf{Attrition Adjustment}: 20\% expected dropout $\rightarrow$ $63 / 0.80 = 79 \approx 80$ total (40 per group)
\end{itemize}

\subsubsection{Missing Data Handling}

\begin{itemize}
    \item \textbf{Primary Approach}: LMM handles missing data under MAR assumption (uses all available data)
    \item \textbf{MAR Assumption Check}: Little's MCAR test ($p > 0.05$ supports MAR); compare baseline characteristics of completers vs non-completers
    \item \textbf{Sensitivity Analysis}: $\delta$-adjustment pattern mixture model (assume missing participants are 0.5 SD worse than completers)
    \item \textbf{Multiple Imputation}: For baseline missing data ($<$5\%), use MICE with m=20 imputations (PMM method)
\end{itemize}

\subsubsection{Software}

All analyses conducted in R (v4.4+):
\begin{itemize}
    \item \textbf{LMM}: \texttt{lme4} package (\texttt{lmer} function)
    \item \textbf{HLM}: \texttt{nlme} package (\texttt{lme} function)
    \item \textbf{Multiple Imputation}: \texttt{mice} package
    \item \textbf{ICC}: \texttt{irr} package (\texttt{icc} function)
    \item \textbf{Effect Sizes}: \texttt{effectsize} package
\end{itemize}

\subsection{Quality Control \& Fidelity}

\textbf{Intervention Fidelity}:
\begin{itemize}
    \item \textbf{Session Recording}: 100\% of sessions logged (AI prompts, participant responses, scores)
    \item \textbf{Manual Audit}: 10\% of sessions randomly selected for manual review
    \item \textbf{Fidelity Criteria}: $\ge$90\% adherence to standardized protocol
    \item \textbf{Corrective Action}: Re-training for interventionists with $<$90\% fidelity
\end{itemize}

\textbf{Assessor Training \& Monitoring}:
\begin{itemize}
    \item \textbf{Initial Training}: 2h video module + 10 practice narratives with gold-standard comparison
    \item \textbf{ICC Calibration}: 20 test narratives; ICC$>$0.90 required for certification
    \item \textbf{Monthly Refresher}: 2h/month to prevent rater drift
    \item \textbf{Ongoing Monitoring}: 10\% of narratives double-scored; ICC tracked monthly; re-training if ICC$<$0.85
\end{itemize}

\textbf{Data Quality}:
\begin{itemize}
    \item \textbf{Source Data Verification}: 10\% of data points verified against source documents
    \item \textbf{Outlier Detection}: Automated flagging of values $\pm$3 SD from mean; clinical review for plausibility
    \item \textbf{Protocol Deviation Log}: All deviations documented with reason and corrective action
\end{itemize}

\section{Ethical Considerations}
\label{sec:ethics}

\subsection{Data Privacy}

Narratives contain sensitive personal information. The scorer is designed for local execution, ensuring data never leaves the user's device. When deployed in research settings, narratives are:
\begin{itemize}
    \item Encrypted at rest (AES-256)
    \item Anonymized (participant IDs, no PII)
    \item Deleted after 5 years (per IRB protocol)
\end{itemize}

\subsection{Emotional Safety}

Discussing personal memories can evoke strong emotions. The scorer includes an emotional arousal detector (not described in this report) that flags high-arousal narratives for human review. Crisis protocols are in place for participants showing signs of distress.

\subsection{Algorithmic Fairness}

The scorer's vocabularies are based on Standard Mandarin. Performance may be lower for:
\begin{itemize}
    \item Dialect speakers (Cantonese, Shanghainese, etc.)
    \item Non-native Chinese speakers
    \item Older adults with atypical speech patterns (e.g., dementia-related)
\end{itemize}

We acknowledge this limitation and plan to expand vocabularies in future versions.

\subsection{Transparency}

As an open-source tool, the scorer's algorithms are fully transparent. Users can inspect, modify, and extend the code. This contrasts with proprietary LLM APIs, where scoring logic is opaque.

\section{Limitations and Future Work}
\label{sec:limitations}

\subsection{Current Limitations}

\begin{enumerate}[label=(\arabic*)]
    \item \textbf{Rule-based extraction}: Does not capture semantic nuance (e.g., implied causality without explicit markers)
    \item \textbf{Simplified Chinese only}: No support for Traditional Chinese or dialects
    \item \textbf{No ASR integration}: Requires text input; speech-to-text not yet integrated
    \item \textbf{Limited validation}: Empirical data pending; current evidence is from mock testing
    \item \textbf{Event boundary detection v2 requires further validation}: The topic-transition-aware algorithm targets F1 $> 0.75$ but has not yet been validated on a large-scale dataset with human-annotated boundaries
\end{enumerate}

\subsection{Future Directions}

\begin{enumerate}[label=(\arabic*)]
    \item \textbf{LLM-enhanced scoring}: Hybrid approach using LLMs for event extraction while retaining interpretable dimension scores
    \item \textbf{Checklist-based LLM scoring}: Decompose subjective dimensions into binary yes/no items following CheckEval \citep{checkeval2024} and Healthcare LLM-Judge \citep{healthcare_llm_judge2025} paradigms, targeting ICC $> 0.80$
    \item \textbf{Cognitive dimension extraction}: Fine-tune Llama-3-8B (QLoRA) to extract 7 cognitive categories (thought, emotion, social interaction, location, time...) following recent work on Parkinson's disease narrative NLP \citep{pd_narrative_nlp2025}
    \item \textbf{Multilingual support}: Extend to Cantonese, English, and other languages spoken by diaspora communities
    \item \textbf{ASR integration}: Whisper or Dolphin \citep{dolphin2025} for end-to-end speech-to-score pipeline, with dialect support (22 Chinese dialects)
    \item \textbf{Longitudinal validation}: 6-month follow-up to assess test-retest reliability and sensitivity to change
    \item \textbf{Clinical utility studies}: Does feedback from the scorer improve RT outcomes?
\end{enumerate}

\section{Conclusion}
\label{sec:conclusion}

The CittaVerse Narrative Scorer v0.6 provides automated, interpretable assessment of Chinese autobiographical memory quality across six dimensions with novel event boundary detection. By combining rule-based feature extraction with LLM-enhanced event segmentation and research-derived scoring heuristics, the tool achieves full automation without requiring cloud-based LLMs for core scoring. The event boundary detection v2 algorithm targets F1 $> 0.75$ against human annotation, extending recent work on LLM-based event segmentation. Open-source release supports community validation and extension. Empirical validation is ongoing through a pilot RCT (N=50, 8-week intervention), with results expected in Q3 2026. We position this work as a foundational contribution to AI-enhanced digital therapeutics for cognitive health in aging populations.

\section*{Acknowledgments}

We thank the CittaVerse team and pilot participants for their contributions. This work was supported by 一念万相科技 (CittaVerse).

\bibliographystyle{apalike}
\bibliography{references}

\appendix

\section{Prompt Templates (for LLM-enhanced version)}
\label{app:prompts}

\subsection*{A.1 Event Boundary Detection Prompt}

\textbf{Role}: Professional narrative analysis assistant.

\textbf{Task}: Analyze the following oral transcript and identify event boundaries.

\begin{verbatim}
You are a professional narrative analysis assistant. Please analyze 
the following oral transcript and identify event boundaries.

[Input Text]
{transcript}

[Tasks]
1. Identify event boundaries (event transition points) in the text
2. Assign a short label to each event
3. Output format: JSON

[Event Boundary Criteria]
- Time changes (e.g., "later", "the next day", "that summer")
- Location changes (e.g., "arrived at school", "returned home")
- Character changes (e.g., "my colleague and I", "later I was alone")
- Theme/activity changes (e.g., "after starting work", "after work")
- Emotional/attitudinal shifts (e.g., "but", "however", "unexpectedly")

[Output Format]
{
  "events": [
    {
      "event_id": 1,
      "start_char": 0,
      "end_char": 150,
      "label": "Childhood memory - first day of school",
      "boundary_cues": ["temporal marker", "location marker"],
      "text_snippet": "I remember that morning..."
    }
  ],
  "total_events": 5,
  "avg_event_length_chars": 120
}
\end{verbatim}

\subsection*{A.2 Six-Dimension Scoring Prompt}

\textbf{Role}: Professional narrative quality assessment expert.

\textbf{Task}: Evaluate the oral transcript using the six-dimension framework.

\begin{verbatim}
Six Dimensions:
1. Specificity: Sensory details, time, place, characters
2. Emotionality: Emotional expression and changes
3. Coherence: Logical connections and causal relationships
4. Reflectiveness: Meaning-making and self-reflection
5. Vividness: Vividness and imagery
6. Completeness: Clear beginning, middle, and end

Scoring: 1-5 scale (1=completely missing, 5=excellent)

Output: JSON with scores, evidence, strengths, and areas for improvement.
\end{verbatim}

\section{Assessment Scales (Chinese translations)}
\label{app:scales}

\subsection*{B.1 System Usability Scale (SUS)}

\textbf{Source}: Brooke (1996). \textit{SUS: A "quick and dirty" usability scale.}

\textbf{Instructions}: Rate each statement based on your actual experience. 1=Strongly Disagree, 5=Strongly Agree.

\begin{center}
\begin{tabular}{clccccc}
\hline
\# & Item & 1 & 2 & 3 & 4 & 5 \\
\hline
1 & I would like to use this system frequently & ○ & ○ & ○ & ○ & ○ \\
2 & I found the system unnecessarily complex & ○ & ○ & ○ & ○ & ○ \\
3 & I thought the system was easy to use & ○ & ○ & ○ & ○ & ○ \\
4 & I think I would need technical support & ○ & ○ & ○ & ○ & ○ \\
5 & I found the various functions well integrated & ○ & ○ & ○ & ○ & ○ \\
6 & I thought there was too much inconsistency & ○ & ○ & ○ & ○ & ○ \\
7 & I would imagine most people learn quickly & ○ & ○ & ○ & ○ & ○ \\
8 & I found the system very cumbersome & ○ & ○ & ○ & ○ & ○ \\
9 & I felt very confident using the system & ○ & ○ & ○ & ○ & ○ \\
10 & I needed to learn a lot before using & ○ & ○ & ○ & ○ & ○ \\
\hline
\end{tabular}
\end{center}

\textbf{Scoring}: 
Odd items: score = value - 1; Even items: score = 5 - value. 
Total = sum × 2.5. Range: 0-100. 
$\geq 68$: above average; $\geq 80.3$: good (A); $\geq 85.5$: excellent (A+).

\subsection*{B.2 Net Promoter Score (NPS)}

\textbf{Question}: How likely are you to recommend this system to a friend or colleague?

\textbf{Scale}: 0 (Not at all likely) to 10 (Extremely likely)

\textbf{Scoring}: 
Detractors (0-6), Passives (7-8), Promoters (9-10). 
NPS = \%Promoters - \%Detractors. Range: -100 to +100.

\subsection*{B.3 Technology Anxiety Scale}

\textbf{Source}: Adapted from Meier (1985).

\textbf{Sample Items} (1=Strongly Disagree, 5=Strongly Agree):
\begin{itemize}
  \item I worry about accidentally breaking the system
  \item Using new systems makes me nervous
  \item I worry about pressing wrong buttons
\end{itemize}

\textbf{Scoring}: Sum of 6 items. Range: 6-30. Higher = more anxiety.

\subsection*{B.4 Privacy Concerns Scale}

\textbf{Source}: Adapted from Iivari \& Iivari (2013).

\textbf{Sample Items} (1=Strongly Disagree, 5=Strongly Agree):
\begin{itemize}
  \item I worry about my personal information being leaked
  \item I worry the system will record my private conversations
  \item I am unsure how my data will be used
\end{itemize}

\textbf{Scoring}: Sum of 6 items. Range: 6-30. Higher = more concern.

\section{Source Code Repository}
\label{app:source-code}

\textbf{GitHub}: \url{https://github.com/cittaverse/narrative-scorer} \\
\textbf{License}: MIT \\
\textbf{Version}: v0.6 \\
\textbf{DOI}: [To be added after arXiv submission]

\textbf{Installation}:
\begin{verbatim}
pip install git+https://github.com/cittaverse/narrative-scorer.git
\end{verbatim}

\textbf{Usage}:
\begin{verbatim}
from narrative_scorer import SixDimensionScorer
scorer = SixDimensionScorer(model="qwen-plus")
result = scorer.score("我记得那天早上...")
\end{verbatim}

\textbf{Citation}:
\begin{verbatim}
@software{narrative_scorer_2026,
  author = {CittaVerse Team},
  title = {CittaVerse Narrative Scorer},
  version = {0.6},
  year = {2026},
  url = {https://github.com/cittaverse/narrative-scorer},
  license = {MIT}
}
\end{verbatim}

\end{document}
